\documentclass[12pt,a4paper]{article}
\usepackage[utf8]{inputenc}
\usepackage[ngerman]{babel}
\usepackage{amsmath}
\usepackage{amssymb}
\usepackage{amsthm}
\usepackage{geometry}
\usepackage{xcolor}
\usepackage{tcolorbox}
\usepackage{enumitem}

\geometry{margin=2.5cm}

\theoremstyle{definition}
\newtheorem{definition}{Definition}
\newtheorem{satz}{Satz}
\newtheorem{beispiel}{Beispiel}

\title{Strategie III: Ski-Rental Problem\\
\large Detaillierte Erklärung}
\author{}
\date{}

\begin{document}

\maketitle

\section{Das Problem}

\begin{tcolorbox}[colback=blue!5!white,colframe=blue!75!black,title=Ski-Rental Problem]
\begin{itemize}
    \item Leihen kostet \textbf{1 Euro/Tag}
    \item Kaufen kostet \textbf{$M$ Euro}
    \item Wir wissen nicht, wie viele Tage $T$ wir Ski brauchen
\end{itemize}
\textbf{Ziel:} Finde eine Strategie, die möglichst nah an der optimalen Lösung ist.
\end{tcolorbox}

\section{Strategie III}

\begin{center}
\colorbox{yellow!50}{\large \textbf{Kaufe am $M$-ten guten Tag}}
\end{center}

\subsection{Was bedeutet das?}

\begin{itemize}[leftmargin=*]
    \item Leihe die Ski für die ersten $M-1$ Tage
    \item Am Tag $M$ kaufst du die Ski
    \item Danach keine weiteren Kosten
\end{itemize}

\section{Analyse des Competitive Ratio}

Die \textbf{optimale Lösung} zahlt:
\[
\text{OPT} = \min(M, T)
\]

\begin{itemize}
    \item Wenn $T < M$: optimal ist nur leihen $\Rightarrow$ OPT $= T$
    \item Wenn $T \geq M$: optimal ist am Tag 1 kaufen $\Rightarrow$ OPT $= M$
\end{itemize}

\subsection{Kosten von Strategie III}

\[
C_{\text{det}} = \underbrace{(M-1)}_{\text{Leihen}} + \underbrace{M}_{\text{Kaufen}} = 2M - 1
\]

\textbf{Erklärung:}
\begin{enumerate}
    \item Wir leihen für $M-1$ Tage $\Rightarrow$ Kosten: $M-1$
    \item Am Tag $M$ kaufen wir $\Rightarrow$ Kosten: $M$
    \item \textbf{Gesamt:} $(M-1) + M = 2M - 1$
\end{enumerate}

\section{Berechnung des Competitive Ratio}

Das Competitive Ratio ist:
\[
\frac{C_{\text{det}}}{\text{OPT}}
\]

Wir müssen zwei Fälle betrachten:

\subsection{Fall 1: $T < M$ (Wenige Tage)}

\begin{tcolorbox}[colback=red!5!white,colframe=red!75!black]
\textbf{Annahme:} $T < M$ (wir brauchen die Ski nur kurz)

\textbf{OPT:} Optimal wäre nur zu leihen $\Rightarrow$ OPT $= T$

\textbf{Strategie III zahlt:} $2M - 1$ (haben schon gekauft, bevor wir aufhören)

\textbf{Competitive Ratio:}
\[
\frac{C_{\text{det}}}{\text{OPT}} = \frac{2M-1}{T}
\]

\textbf{Worst-Case:} Wenn $T$ so klein wie möglich ist (aber $T \in \mathbb{N}^+$), ist das Ratio am größten.

Sei $T = X$ mit $X < M$. Dann:
\[
\frac{C_{\text{det}}}{\text{OPT}} = \frac{2M-1}{X} \geq \frac{2M-1}{\min(X,M)} = \frac{2M-1}{X}
\]

Um eine obere Schranke zu finden, betrachten wir:
\[
\frac{X-1+M}{X} = \frac{X-1}{X} + \frac{M}{X} = 1 - \frac{1}{X} + \frac{M}{X}
\]

Für $X < M$ gilt: Die Strategie zahlt $(X-1)$ fürs Leihen + $M$ fürs Kaufen, aber OPT = $X$.

\textbf{Alternative Betrachtung:}
\[
\frac{X-1+X+1}{X} = \frac{2X}{X} = 2
\]

Hier wird approximiert: Im worst-case (direkt vor dem Kaufen aufhören) zahlen wir ca. doppelt so viel wie OPT.
\end{tcolorbox}

\subsection{Fall 2: $T \geq M$ (Viele Tage)}

\begin{tcolorbox}[colback=green!5!white,colframe=green!75!black]
\textbf{Annahme:} $T \geq M$ (wir brauchen die Ski lange)

\textbf{OPT:} Optimal ist am Tag 1 zu kaufen $\Rightarrow$ OPT $= M$

\textbf{Strategie III zahlt:} $2M - 1$

\textbf{Competitive Ratio:}
\[
\frac{C_{\text{det}}}{\text{OPT}} = \frac{2M-1}{M} = \frac{2M}{M} - \frac{1}{M} = 2 - \frac{1}{M}
\]

\textbf{Interpretation:}
\begin{itemize}
    \item Wir zahlen fast doppelt so viel wie optimal
    \item Für großes $M$: Ratio $\approx 2$
    \item Für $M \to \infty$: Ratio $\to 2$
\end{itemize}
\end{tcolorbox}

\section{Die Formel im Bild erklärt}

Im Bild steht:
\[
\frac{C_{\text{det}}}{C_{\text{OPT}}} = \frac{X-1+M}{\min(X,M)} \geq \min\left(\frac{X-1+X+1}{X}, \frac{M-1+M}{M}\right)
\]

\subsection{Bedeutung der Terme}

\begin{enumerate}
    \item \textbf{$\displaystyle\frac{X-1+X+1}{X}$} (Fall $X < M$):
    \begin{itemize}
        \item Zähler: $(X-1)$ Leihen + $(X+1)$ annähernde Kaufkosten
        \item Vereinfachung: $\displaystyle\frac{2X}{X} = 2$
        \item Dies ist der worst-case, wenn man kurz vor dem Kaufzeitpunkt aufhört
    \end{itemize}
    
    \item \textbf{$\displaystyle\frac{M-1+M}{M}$} (Fall $X \geq M$):
    \begin{itemize}
        \item Zähler: $(M-1)$ Tage Leihen + $M$ Kaufpreis = $2M-1$
        \item Nenner: OPT $= M$ (hätte sofort kaufen sollen)
        \item Ergebnis: $\displaystyle 2 - \frac{1}{M}$
    \end{itemize}
\end{enumerate}

\subsection{Das Minimum}

\[
\min\left(2, 2-\frac{1}{M}\right) = 2 - \frac{1}{M} < 2
\]

$\Rightarrow$ Das Competitive Ratio ist \textbf{höchstens} $2 - \frac{1}{M}$ für endliches $M$.

Für $M \to \infty$: Competitive Ratio $\to 2$

\section{Hauptsatz}

\begin{satz}
Strategie III ist \textbf{deterministisch} und \textbf{2-kompetitiv}.

Keine deterministische Strategie kann besser als $2-\frac{1}{M}$-kompetitiv sein (für $M \to \infty$: nicht besser als 2-kompetitiv).
\end{satz}

\section{Beispiele}

\begin{beispiel}[Kurze Nutzung: $T = 3$, $M = 10$]
\begin{itemize}
    \item \textbf{OPT:} Nur leihen $\Rightarrow$ Kosten $= 3$
    \item \textbf{Strategie III:} Leihe 9 Tage, kaufe am Tag 10 $\Rightarrow$ Kosten $= 9 + 10 = 19$
    \item \textbf{Ratio:} $\displaystyle\frac{19}{3} \approx 6.33$ (schlecht, weil $T \ll M$)
\end{itemize}
\end{beispiel}

\begin{beispiel}[Mittlere Nutzung: $T = 10$, $M = 10$]
\begin{itemize}
    \item \textbf{OPT:} Sofort kaufen $\Rightarrow$ Kosten $= 10$
    \item \textbf{Strategie III:} Leihe 9 Tage, kaufe am Tag 10 $\Rightarrow$ Kosten $= 9 + 10 = 19$
    \item \textbf{Ratio:} $\displaystyle\frac{19}{10} = 1.9$
\end{itemize}
\end{beispiel}

\begin{beispiel}[Lange Nutzung: $T = 100$, $M = 10$]
\begin{itemize}
    \item \textbf{OPT:} Sofort kaufen $\Rightarrow$ Kosten $= 10$
    \item \textbf{Strategie III:} Leihe 9 Tage, kaufe am Tag 10 $\Rightarrow$ Kosten $= 9 + 10 = 19$
    \item \textbf{Ratio:} $\displaystyle\frac{19}{10} = 1.9$
\end{itemize}
\end{beispiel}

\begin{beispiel}[Großes $M$: $M = 1000$]
\begin{itemize}
    \item \textbf{Kosten Strategie III:} $2 \cdot 1000 - 1 = 1999$
    \item \textbf{OPT (lange Nutzung):} $1000$
    \item \textbf{Ratio:} $\displaystyle\frac{1999}{1000} = 1.999 \approx 2$
\end{itemize}
\end{beispiel}

\section{Strategie IV: Randomisierte Strategie}

\begin{center}
\colorbox{orange!50}{\large \textbf{Können wir die 2-kompetitive Schranke brechen? Mit Randomisierung!}}
\end{center}

\subsection{Die Idee}

Strategie III ist deterministisch und erreicht ein Competitive Ratio von $2 - \frac{1}{M} \approx 2$. Keine deterministische Strategie kann besser sein. Aber was ist mit \textbf{randomisierten} Strategien?

\begin{tcolorbox}[colback=orange!5!white,colframe=orange!75!black,title=Strategie IV]
\textbf{Werfe eine Münze:}
\begin{itemize}
    \item \textbf{Kopf:} Kaufe am $M$-ten guten Tag
    \item \textbf{Zahl:} Kaufe am $\alpha M$-ten guten Tag (mit $\alpha \in (0,1)$)
\end{itemize}

Jede Option hat Wahrscheinlichkeit $\frac{1}{2}$.
\end{tcolorbox}

\subsection{Parameter $\alpha$}

Der Parameter $\alpha$ bestimmt, wann wir im zweiten Fall kaufen:
\begin{itemize}
    \item Wenn $\alpha = \frac{1}{2}$: Kaufe am $\frac{M}{2}$-ten Tag (die Hälfte)
    \item Wenn $\alpha$ kleiner ist: Kaufe früher
    \item Wenn $\alpha$ größer ist: Kaufe später
\end{itemize}

Wir werden sehen, dass $\alpha = \frac{1}{2}$ eine gute Wahl ist!

\subsection{Erwartete Kosten}

Da die Strategie randomisiert ist, berechnen wir die \textbf{erwarteten Kosten}:

\[
E[c_{\text{Strategie IV}}] = \frac{1}{2} \cdot c_{\text{Kopf}} + \frac{1}{2} \cdot c_{\text{Zahl}}
\]

\textbf{Kosten bei Kopf (kaufe am Tag $M$):}
\[
c_{\text{Kopf}} = (M-1) + M = 2M - 1
\]

\textbf{Kosten bei Zahl (kaufe am Tag $\alpha M$):}
\[
c_{\text{Zahl}} = (\alpha M - 1) + M = (1+\alpha)M - 1
\]

\subsection{Analyse: Fall $T = M$}

\textbf{Beobachtung:} Der schlechteste Fall für diese Strategie ist, wenn $T = M$ oder $T = \alpha M$.

\begin{tcolorbox}[colback=purple!5!white,colframe=purple!75!black,title=Fall $T = M$]

\textbf{OPT:} Optimal ist am Tag 1 kaufen $\Rightarrow$ OPT $= M$

\textbf{Erwartete Kosten:}
\begin{align*}
E[c_{\text{Strategie IV}}] &= \frac{1}{2} \cdot (2M-1) + \frac{1}{2} \cdot ((1+\alpha)M-1) \\
&= \frac{1}{2}(2M - 1 + (1+\alpha)M - 1) \\
&= \frac{1}{2}(2M + (1+\alpha)M - 2) \\
&= \frac{1}{2}(3M + \alpha M - 2) \\
&= \frac{1}{2}((3+\alpha)M - 2) \\
&= \frac{(3+\alpha)M - 2}{2}
\end{align*}

\textbf{Competitive Ratio:}
\begin{align*}
\frac{E[c_{\text{Strategie IV}}]}{c_{\text{OPT}}} &= \frac{\frac{1}{2}(2M-1) + \frac{1}{2}((1+\alpha)M-1)}{M} \\
&= \frac{(3+\alpha)M - 2}{2M} \\
&= \frac{3+\alpha}{2} - \frac{1}{M} \\
&\xrightarrow{M \to \infty} \frac{3+\alpha}{2}
\end{align*}

\textbf{Interpretation:}
\begin{itemize}
    \item Das Ratio hängt von $\alpha$ ab!
    \item Für $\alpha = \frac{1}{2}$: Ratio $= \frac{3+0.5}{2} = \frac{3.5}{2} = 1.75$
    \item Dies ist besser als 2!
\end{itemize}
\end{tcolorbox}

\subsection{Analyse: Fall $T = \alpha M$}

\begin{tcolorbox}[colback=cyan!5!white,colframe=cyan!75!black,title=Fall $T = \alpha M$]

\textbf{OPT:} Optimal ist am Tag 1 kaufen $\Rightarrow$ OPT $= M$

Moment! Wenn $T = \alpha M < M$, dann ist optimal nur zu leihen:
\[
\text{OPT} = \alpha M
\]

\textbf{Erwartete Kosten:}
\begin{itemize}
    \item \textbf{Kopf} (kaufe am Tag $M$): Wir haben schon bei $T = \alpha M$ aufgehört, also: $c = \alpha M$ (nur Leihkosten)
    \item \textbf{Zahl} (kaufe am Tag $\alpha M$): $c = (\alpha M - 1) + M = (1+\alpha)M - 1$
\end{itemize}

\begin{align*}
E[c_{\text{Strategie IV}}] &= \frac{1}{2} \cdot \alpha M + \frac{1}{2} \cdot ((1+\alpha)M - 1) \\
&= \frac{1}{2}(\alpha M + (1+\alpha)M - 1) \\
&= \frac{1}{2}(2\alpha M + M - 1) \\
&= \frac{(2\alpha + 1)M - 1}{2}
\end{align*}

\textbf{Competitive Ratio:}
\begin{align*}
\frac{E[c_{\text{Strategie IV}}]}{c_{\text{OPT}}} &= \frac{\frac{1}{2} \cdot \alpha M + \frac{1}{2} \cdot ((1+\alpha)M - 1)}{\alpha M} \\
&= \frac{(2\alpha + 1)M - 1}{2\alpha M} \\
&= \frac{2\alpha + 1}{2\alpha} - \frac{1}{2\alpha M} \\
&= \frac{2\alpha}{2\alpha} + \frac{1}{2\alpha} - \frac{1}{2\alpha M} \\
&= 1 + \frac{1}{2\alpha} - \frac{1}{2\alpha M} \\
&\xrightarrow{M \to \infty} 1 + \frac{1}{2\alpha}
\end{align*}

\textbf{Interpretation:}
\begin{itemize}
    \item Das Ratio hängt von $\alpha$ ab!
    \item Je kleiner $\alpha$, desto größer das Ratio!
    \item Für $\alpha = \frac{1}{2}$: Ratio $= 1 + \frac{1}{2 \cdot 0.5} = 1 + 1 = 2$
\end{itemize}
\end{tcolorbox}

\subsection{Optimale Wahl von $\alpha$}

Wir haben zwei kritische Fälle:
\begin{enumerate}
    \item Fall $T = M$: Ratio $= \frac{3+\alpha}{2}$ (für $M \to \infty$)
    \item Fall $T = \alpha M$: Ratio $= 1 + \frac{1}{2\alpha}$ (für $M \to \infty$)
\end{enumerate}

Um das \textbf{maximale Ratio zu minimieren}, setzen wir die beiden gleich:

\[
\frac{3+\alpha}{2} = 1 + \frac{1}{2\alpha}
\]

Lösen nach $\alpha$:
\begin{align*}
\frac{3+\alpha}{2} &= 1 + \frac{1}{2\alpha} \\
3 + \alpha &= 2 + \frac{1}{\alpha} \\
\alpha + 1 &= \frac{1}{\alpha} \\
\alpha^2 + \alpha &= 1 \\
\alpha^2 + \alpha - 1 &= 0
\end{align*}

Mit der quadratischen Formel:
\[
\alpha = \frac{-1 \pm \sqrt{1 + 4}}{2} = \frac{-1 \pm \sqrt{5}}{2}
\]

Da $\alpha \in (0,1)$:
\[
\alpha = \frac{-1 + \sqrt{5}}{2} \approx 0.618
\]

Dies ist der \textbf{goldene Schnitt} minus 1!

\textbf{Competitive Ratio mit optimalem $\alpha$:}
\[
\frac{3+\alpha}{2} = \frac{3 + \frac{\sqrt{5}-1}{2}}{2} = \frac{\frac{6 + \sqrt{5} - 1}{2}}{2} = \frac{5 + \sqrt{5}}{4} \approx 1.809
\]

\subsection{Alternative: $\alpha = \frac{1}{2}$}

Für einfachere Berechnungen kann man auch $\alpha = \frac{1}{2}$ wählen:

\textbf{Fall $T = M$:}
\[
\text{Ratio} = \frac{3 + 0.5}{2} = 1.75
\]

\textbf{Fall $T = \alpha M = \frac{M}{2}$:}
\[
\text{Ratio} = 1 + \frac{1}{2 \cdot 0.5} = 1 + 1 = 2
\]

\textbf{Maximum:} $\max(1.75, 2) = 2$

Also ist mit $\alpha = \frac{1}{2}$ die Strategie immer noch 2-kompetitiv, aber im Fall $T = M$ nur 1.75-kompetitiv!

\subsection{Erklärung der Formeln aus dem Bild}

\begin{enumerate}
    \item \textbf{Fall $T = M$:}
    \[
    \frac{E[c_{\text{Strategie IV}}]}{c_{\text{OPT}}} = \frac{\frac{1}{2} \cdot (2M-1) + \frac{1}{2} \cdot ((1+\alpha)M-1)}{M} = \frac{3+\alpha}{2} - \frac{1}{M} \xrightarrow{M \to \infty} \frac{3+\alpha}{2}
    \]
    
    \textbf{Woher kommt $(3+\alpha)$?}
    \begin{align*}
    \text{Zähler} &= \frac{1}{2}(2M - 1) + \frac{1}{2}((1+\alpha)M - 1) \\
    &= \frac{1}{2}(2M + (1+\alpha)M - 2) \\
    &= \frac{1}{2}((3+\alpha)M - 2)
    \end{align*}
    
    Dividiert durch $M$:
    \[
    \frac{(3+\alpha)M - 2}{2M} = \frac{3+\alpha}{2} - \frac{1}{M}
    \]
    
    \item \textbf{Fall $T = \alpha M$:}
    \[
    \frac{E[c_{\text{Strategie IV}}]}{c_{\text{OPT}}} = \frac{\frac{1}{2} \cdot \alpha M + \frac{1}{2} \cdot ((1+\alpha)M-1)}{\alpha M} = 1 + \frac{1}{2\alpha} - \frac{1}{2\alpha M} \xrightarrow{M \to \infty} 1 + \frac{1}{2\alpha}
    \]
    
    \textbf{Woher kommt $1 + \frac{1}{2\alpha}$?}
    \begin{align*}
    \text{Zähler} &= \frac{1}{2} \alpha M + \frac{1}{2}((1+\alpha)M - 1) \\
    &= \frac{1}{2}(\alpha M + (1+\alpha)M - 1) \\
    &= \frac{1}{2}((2\alpha + 1)M - 1)
    \end{align*}
    
    Dividiert durch $\alpha M$:
    \begin{align*}
    \frac{(2\alpha + 1)M - 1}{2\alpha M} &= \frac{(2\alpha + 1)M}{2\alpha M} - \frac{1}{2\alpha M} \\
    &= \frac{2\alpha + 1}{2\alpha} - \frac{1}{2\alpha M} \\
    &= \frac{2\alpha}{2\alpha} + \frac{1}{2\alpha} - \frac{1}{2\alpha M} \\
    &= 1 + \frac{1}{2\alpha} - \frac{1}{2\alpha M}
    \end{align*}
\end{enumerate}

\subsection{Beispiel: $\alpha = \frac{1}{2}$, $M = 100$}

\begin{beispiel}[Randomisierte Strategie mit $\alpha = 0.5$, $M = 100$]

\textbf{Fall 1: $T = M = 100$}
\begin{itemize}
    \item \textbf{OPT:} Kaufe sofort $\Rightarrow$ Kosten $= 100$
    \item \textbf{Strategie IV:}
    \begin{itemize}
        \item Kopf (50\%): Leihe 99 Tage, kaufe am Tag 100 $\Rightarrow$ $99 + 100 = 199$
        \item Zahl (50\%): Leihe 49 Tage, kaufe am Tag 50 $\Rightarrow$ $49 + 100 = 149$
        \item Erwartete Kosten: $\frac{1}{2}(199) + \frac{1}{2}(149) = 174$
    \end{itemize}
    \item \textbf{Ratio:} $\frac{174}{100} = 1.74$ (besser als 2!)
\end{itemize}

\textbf{Fall 2: $T = \alpha M = 50$}
\begin{itemize}
    \item \textbf{OPT:} Nur leihen $\Rightarrow$ Kosten $= 50$
    \item \textbf{Strategie IV:}
    \begin{itemize}
        \item Kopf (50\%): Nur leihen (haben noch nicht gekauft) $\Rightarrow$ $50$
        \item Zahl (50\%): Leihe 49 Tage, kaufe am Tag 50 $\Rightarrow$ $49 + 100 = 149$
        \item Erwartete Kosten: $\frac{1}{2}(50) + \frac{1}{2}(149) = 99.5$
    \end{itemize}
    \item \textbf{Ratio:} $\frac{99.5}{50} = 1.99$ (fast 2)
\end{itemize}
\end{beispiel}

\section{Zusammenfassung}

\subsection{Strategie III (Deterministisch)}
\begin{itemize}
    \item \textbf{Strategie:} Kaufe am $M$-ten Tag
    \item \textbf{Kosten:} $2M - 1$ (unabhängig von $T$, solange $T \geq M$)
    \item \textbf{Competitive Ratio:} Maximal $2 - \frac{1}{M}$ $\approx$ 2
    \item \textbf{Optimal:} Keine deterministische Strategie kann besser sein
\end{itemize}

\subsection{Strategie IV (Randomisiert)}
\begin{itemize}
    \item \textbf{Strategie:} Mit Wahrscheinlichkeit $\frac{1}{2}$ kaufe am $M$-ten Tag, sonst am $\alpha M$-ten Tag
    \item \textbf{Erwartete Kosten:} Hängt von $\alpha$ ab
    \item \textbf{Competitive Ratio:} Mit optimalem $\alpha \approx 0.618$: ca. 1.809
    \item \textbf{Mit $\alpha = \frac{1}{2}$:} Ratio zwischen 1.75 und 2
    \item \textbf{Vorteil:} Randomisierung durchbricht die 2-kompetitive Schranke!
\end{itemize}

\subsection{Vergleich}

\begin{center}
\begin{tabular}{|l|c|c|}
\hline
\textbf{Strategie} & \textbf{Typ} & \textbf{Comp. Ratio} \\
\hline
Strategie III & Deterministisch & $\approx 2$ \\
Strategie IV ($\alpha = 0.5$) & Randomisiert & $\leq 2$ (oft $\approx 1.75$) \\
Strategie IV ($\alpha \approx 0.618$) & Randomisiert & $\approx 1.809$ \\
\hline
\end{tabular}
\end{center}

\textbf{Fazit:} Randomisierung kann deterministische untere Schranken durchbrechen!

\end{document}
